% !TEX root = ../main.tex
% capitulos/2_metodologia.tex

\subsection{Descrição}

Aplica-se nos casos em que o trabalho segue uma linha de desenvolvimento de assuntos contínuos, conforme a estrutura abaixo sugerida:

Apresentar o assunto estudado, abordando os aspectos gerais e buscando introduzir o leitor na temática delineada. Também, fazer uma descrição sucinta dos objetivos da pesquisa. Ressaltar a importância da pesquisa dentro do contexto científico e/ou tecnológico, relatando as possíveis contribuições dos resultados alcançados.

Abordar os aspectos teóricos diretamente relacionados com o trabalho desenvolvido, detalhando os assuntos principais do estudo em questão e baseando-se nas diferentes abordagens pesquisadas na literatura\footnote{Literatura atual sobre o assunto (exemplo de como inserir a nota de rodapé).} (livros, teses, dissertações, artigos, trabalhos de congresso, etc.).

Apresentar os materiais e equipamentos utilizados na pesquisa de campo e/ou experimental, detalhando a metodologia e procedimentos empregados durante as atividades para a resolução do problema, os equipamentos e softwares usados no estudo.

Apresentar os resultados alcançados na pesquisa, analisando e discutindo os diversos aspectos de interesse.

Relacionar as conclusões ou considerações finais obtidas de acordo com os resultados observados na pesquisa, podendo incluir sugestões para trabalhos futuros.

Relacionar toda referência citada no artigo. Veja exemplos na seção Referências.
